%%==================================================================================
%%  اللغة الرياضية العربية: مُجمّع AOT ذاتي الاستضافة
%%  Arabic Mathematical Language: A Self-Hosting AOT Compiler with Linear
%%  Ownership and Algebraic Data Types
%%==================================================================================
\documentclass[sigconf]{acmart}

\AtBeginDocument{%
  \providecommand\BibTeX{{%
    Bib\TeX}}}

\setcopyright{acmcopyright}
\copyrightyear{2026}
\acmYear{2027}
\acmConference[PLDI '27]{Programming Language Design and Implementation}{June 2027}{Copenhagen, Denmark}
\acmBooktitle{PLDI '27: Proceedings of the 48th ACM SIGPLAN Conference on Programming Language Design and Implementation}
\acmDOI{10.1145/XXXXXX.XXXXXX}
\acmISBN{978-1-4503-XXXX-X/27/06}

\usepackage{polyglossia}
\usepackage{fontspec}
\usepackage{booktabs}
\usepackage{multirow}
\usepackage{graphicx}
\usepackage{listings}
\usepackage{tikz}
\usepackage{amsmath}
\usepackage{hyperref}

\setdefaultlanguage{arabic}
\setotherlanguage{english}
\newfontfamily\arabicfont[Script=Arabic]{Amiri}
\newfontfamily\arabicfontsf[Script=Arabic]{Amiri}
\newfontfamily\arabicfonttt[Script=Arabic,Scale=0.9]{Amiri}
\setmainfont{Amiri}
\setsansfont{Noto Sans Arabic}

\lstdefinestyle{arlang}{
  basicstyle=\ttfamily\small,
  numbers=none,
  breaklines=true,
  frame=single,
  language=Python,
}

\begin{document}

\title[اللغة الرياضية العربية]{اللغة الرياضية العربية: مُجمّع AOT ذاتي الاستضافة\\مع ملكية خطية وأنواع جبرية}
\subtitle{Arabic Mathematical Language: A Self-Hosting AOT Compiler with Linear Ownership and Algebraic Data Types}

\author{Hamdani Sidi Mohamed}
\authornote{حمداني سيدي محمد --- Hamdani Sidi Mohamed}
\affiliation{%
  \institution{Independent Researcher}
  \city{Morocco}
  \country{Morocco}
}
\email{elhamdolillah@example.com}

\begin{abstract}
We present the Arabic Mathematical Language, the first complete Arabic programming
language that compiles Ahead-of-Time (AOT) to standalone ELF binaries fully
independent of the C standard library (libc). The compiler, implemented in Python,
generates x86\_64 assembly through NASM and integrates modern programming language
techniques: \textit{linear ownership} inspired by Rust for memory safety at compile
time, \textit{pattern matching} in the style of Haskell and Erlang,
\textit{algebraic data types} in the style of OCaml and Rust, and
\textit{traits/typeclasses} in the style of Rust and Haskell. The compiler supports
true parallelism through raw Linux syscalls (\texttt{clone}, \texttt{mmap},
\texttt{wait4}), channels through \texttt{pipe2}, and result/option types with a
dedicated \texttt{?} early-extraction operator. A 91-test suite across 8 harnesses
passes in full, and the compiler demonstrates partial self-hosting capability.
We discuss the language's mathematical constitution (8 principles), the technical
architecture, and the experimental evaluation.
\end{abstract}

\begin{CCSXML}
<ccs2012>
<concept>
<concept_id>10011007.10010940.10010941.10010943</concept_id>
<concept_name>Software and its engineering~Compilers</concept_name>
<concept_significance>500</concept_significance>
</concept>
<concept>
<concept_id>10011007.10011006.10011071</concept_id>
<concept_name>Software and its engineering~Programming languages</concept_name>
<concept_significance>500</concept_significance>
</concept>
<concept>
<concept_id>10003752.10003778.10010124</concept_id>
<concept_name>Theory of computation~Type theory</concept_name>
<concept_significance>300</concept_significance>
</concept>
</ccs2012>
\end{CCSXML}

\ccsdesc[500]{Software and its engineering~Compilers}
\ccsdesc[500]{Software and its engineering~Programming languages}
\ccsdesc[300]{Theory of computation~Type theory}

\keywords{Arabic programming language, AOT compiler, linear ownership, algebraic
data types, traits, self-hosting, x86\_64}

\maketitle

\section{مقدمة}
\label{sec:intro}

\subsection{الدوافع}
\label{sec:intro:motivation}

اللغة العربية، بلغة أكثر من 400 مليون متحدث، تفتقر إلى لغة برمجة كاملة حديثة
تُعبّر عن مفاهيم علوم الحاسب بصيغتها الأصلية. المحاولات السابقة (مثل «ضاد»
و«كلك») ركزت على الترجمة الحرفية للغات موجودة دون تصميم لغوي مستقل، واعتمدت على
مُفسّرات تُشغَّل فوق \texttt{runtime} جاهز. هدفنا بناء لغة عربية \emph{أصلية}
تستوعب أحسن ما في اللغات العالمية وتُعيد إنتاجه بصيغة عربية رياضية، من التحليل
المعجمي حتى توليد ثنائي ELF مستقل تماماً عن مكتبة C القياسية.

\subsection{المشكلة}
\label{sec:intro:problem}

لا توجد لغة برمجة عربية تجمع بين أربعة متطلبات أساسية: (1) \textbf{الأمان} عبر
ضمان أمان الذاكرة في وقت التجميع، (2) \textbf{الأداء} عبر تجميع AOT إلى كود
أصلي بدون runtime ولا جامع قمامة، (3) \textbf{التعبيرية} عبر أنواع جبرية وتطابق
أنماطي وسمات، (4) \textbf{الاستقلالية} عبر عدم الاعتماد على مكتبات خارجية.

\subsection{الحل}
\label{sec:intro:solution}

المُجمّع ذو المراحل التالية:

\begin{enumerate}
\item محلل معجمي (Python) ينتج رموزاً مع تتبع السطر والعمود.
\item محلل نحوي ينتج شجرة بناء مجردة (AST).
\item مدقق ملكية خطية يرفض use-after-move في وقت التجميع.
\item استنتاج أنواع بسيط.
\item مولّد كود ينتج x86\_64 Assembly.
\item تجميع عبر NASM وربط عبر \texttt{ld} إلى ELF بلا libc، يستخدم
  \texttt{sys\_write} و\texttt{sys\_exit} و\texttt{arena\_alloc} المبنية على
  \texttt{mmap}.
\end{enumerate}

\subsection{المساهمات}
\label{sec:intro:contributions}

المساهمات الرئيسية لهذه الورقة هي: (1) أول لغة برمجة عربية كاملة تُجمَّع AOT إلى
ثنائيات ELF مستقلة؛ (2) دستور رياضي من ثمانية مبادئ يحكم كل قرار تصميمي؛ (3)
إعادة إنتاج ذاتية لأحدث تقنيات لغات البرمجة (Rust وHaskell وOCaml) بصيغة عربية؛
(4) مُجمّع ذاتي الاستضافة جزئياً؛ (5) مجموعة اختبار من 91 حالة ناجحة عبر 8
سلاسل.

\section{الخلفية والعمل السابق}
\label{sec:background}

\subsection{لغات البرمجة العربية}
\label{sec:background:arabic}

المحاولات المبكرة لم تتجاوز ثلاثة أنماط: الترجمة الحرفية لـ Python (ضاد)،
الآلات الحاسبة بواجهة عربية (كلك)، واللغات التعليمية المبسطة (جيم). القيود
المشتركة هي الاعتماد الكامل على runtime، وغياب نظام أنواع متقدم، وغياب مفاهيم
ملكية الذاكرة والتطابق الأنماطي.

\subsection{المُجمّعات AOT}
\label{sec:background:aot}

النهج التقليدي (GCC وClang وRustc) يولّد كوداً أصلياً عبر سلاسل معقدة تعتمد على
مكتبات النظام. نهجنا أبسط وأكثر استقلالية: Python تنتج x86\_64 Assembly يُجمَّع
عبر NASM ويُربط عبر \texttt{ld}، دون LLVM ودون libc ودون GC؛ الاعتماديات الوحيدة
هي NASM و\texttt{ld} المتاحان في كل توزيع Linux.

\subsection{نظرية الأنواع الحديثة}
\label{sec:background:types}

نستند إلى ثلاثة أسس نظرية: الأنواع الخطية (Linear Types) كما صاغها Girard
\cite{girard1987} التي تضمن أن كل مورد يُستخدم مرة واحدة بالضبط، وهي أساس
ملكية Rust؛ والأنواع الجبرية (Algebraic Data Types) كما في أعمال McCarthy
\cite{mccarthy1963} التي تجمع بين sum types وproduct types؛ والطبقات (Typeclasses)
كما حددها Wadler وBlott \cite{wadler1989} التي تؤسس للتعدد الصوري ad-hoc
polymorphism.

\section{الدستور الرياضي}
\label{sec:constitution}

\subsection{المبادئ الثمانية}
\label{sec:constitution:principles}

صممنا الدستور كإطار رياضي صارم يحكم كل قرار تصميمي في اللغة:

\begin{table}[h]
\centering
\caption{المبادئ الثمانية للدستور}
\label{tab:principles}
\begin{tabular}{cllc}
\toprule
\textbf{\#} & \textbf{المبدأ} & \textbf{التعريف الرياضي} & \textbf{التطبيق} \\
\midrule
1 & الإحكام (Precision) & $\forall s \in \Sigma:\ |\mathrm{interpret}(s)| = 1$ & لكل رمز معنى واحد \\
2 & التيسير (Simplicity) & $|\Sigma| = 20$ & عشرون رمزاً Turing-complete \\
3 & العدل (Balance) & $\forall f:\ \mathrm{cost}(f) > 0 \land \mathrm{benefit}(f) > \mathrm{cost}(f)$ & لا ميزة بلا ثمن \\
4 & الأمانة (Ownership) & $\forall r,z:\ \exists! o:\ \mathrm{owns}(o,r,z)$ & مالك واحد لكل مورد \\
5 & البيان (Transparency) & $\forall o:\ \mathrm{documented}(o) \land \mathrm{inspectable}(o)$ & كل عملية موثقة \\
6 & الحفظ (Safety) & $\forall p:\ \mathrm{compiled}(p) \Rightarrow \mathrm{safe}(\mathrm{run}(p))$ & أمان بالبناء \\
7 & التفكر (Verifiability) & $\forall e:\ \mathrm{avoidable}(e) \Rightarrow \mathrm{caught\_at\_compile}(e)$ & الأخطاء تُكتشف مبكراً \\
8 & الشمولية (Universality) & أخذ أحسن ما في العالم وإعادة إنتاجه بالعربية & استيعاب + إعادة إنتاج \\
\bottomrule
\end{tabular}
\end{table}

\subsection{التقنيات المُدمَجة والمتجاوزة}
\label{sec:constitution:tech}

أعدنا إنتاج: الملكية الخطية (Rust) كـ\ \texttt{⊸ انقل}، والتطابق الأنماطي
(Erlang/Rust) كـ\ \texttt{طابق}، والأنواع الجبرية (Haskell/OCaml) كـ\ \texttt{نوع}،
والسمات (Rust/Haskell) كـ\ \texttt{سمة}، والقنوات (Go) عبر \texttt{pipe2}،
والخيوط الأصلية عبر \texttt{clone}. وتجاوزنا: جامع القمامة لصالح \texttt{Arena +
Linear Ownership}، وnull لصالح Result/Option، والكائنات لصالح السمات،
والاستثناءات لصالح Result types، وانعكاس runtime لصالح Macros (مرحلة 48).

\section{التصميم والتنفيذ}
\label{sec:design}

\subsection{المحلل المعجمي}
\label{sec:design:lexer}

كل رمز يحمل \texttt{(نوع، قيمة، سطر، عمود)} مما يتيح رسائل أخطاء دقيقة. تدعم
القاموس 20 رمزاً أساسياً، وأسماء بديلة عربية كاملة مثل
\texttt{"دالة":"λ"} و\texttt{"اطبع":"⎕"} و\texttt{"طالما":"μ"} و\texttt{"لكل":"∀"}،
كما يتجاهل المحلل الحركات العربية (التشكيل) في المعرفات.

\subsection{المحلل النحوي وAST}
\label{sec:design:parser}

بنية AST موحدة: التعابير تُمثَّل tuples مثل \texttt{("عدد", 42)} و\texttt{("استدعاء",
"دالة", [args])}، والبيانات مثل \texttt{("أسند", "var", expr)} و\texttt{("لكل",
"var", list, body)}. يدعم المحلل سبعة أنماط: الحرفي، والمتغير، والشامل (\texttt{\_})،
والقائمة (\texttt{⟨أ، ...ذيل⟩})، والشرطي (\texttt{حيث})، والبديل (\texttt{|})،
والباني (\texttt{دائرة(ن)}).

\subsection{الملكية الخطية}
\label{sec:design:ownership}

المدقق يطبق المبدأ الرياضي $\forall r,z:\ \exists! o:\ \mathrm{owns}(o,r,z)$:
عند النقل \texttt{ب ⊸ أ} تُسجَّل \texttt{أ} كمستهلكة، وأي استخدام لاحق يُنتج
خطأ تجميع \emph{«مستهلك — لا يمكن استخدامه»}، مما يحقق مبدأ الحفظ (المبدأ 6).

\subsection{توليد الكود}
\label{sec:design:codegen}

يولّد \texttt{compile\_expr} تعليمات x86\_64 كسلسلة نصوص. Layout الـ heap موحد
لجميع البنيات:

\begin{verbatim}
[ptr + 0]  = len / عدد الحقول
[ptr + 8]  = tag (الوسم) — 0 للقوائم
[ptr+16+k*8] = الحقل k / العنصر k
\end{verbatim}

\textbf{Arena Allocator}: تُبنى الدالة \texttt{arena\_alloc} على \texttt{mmap}
بمحاذاة 16 بايت، ولا يوجد جامع قمامة أبداً. \textbf{فحص الحدود}: كل فهرسة قائمة
تفحص الفهرس ضد الطول وتُنهى البرنامج برسالة خطأ عربية بدل القراءة خارج النطاق.

\subsection{التوازي والقنوات}
\label{sec:design:concurrency}

التوازي الفعلي عبر syscalls خام: \texttt{clone} (56) لإنشاء خيوط بمشاركة
\texttt{CLONE\_VM}، و\texttt{mmap} (9) لمكدس مستقل 64KB، و\texttt{wait4} (61)
للمزامنة. القنوات عبر \texttt{pipe2} (293): \texttt{قناة()} يعيد بنية تحوي
\texttt{read\_fd} و\texttt{write\_fd}، و\texttt{أرسل}/\texttt{استقبل} تستخدم
\texttt{sys\_write}/\texttt{sys\_read} مباشرة.

\subsection{التطابق الأنماطي والأنواع الجبرية}
\label{sec:design:pattern}

\texttt{طابق قيمة : ﴿ نمط ⇒ تعبير ⋄ ... ﴾} يُولَّد كسلسلة \texttt{cmp/jmp} لكل
نمط باني، مع ربط متغيرات النمط في بيئة محلية. الأنواع الجبرية:

\begin{lstlisting}[style=arlang]
نوع شكل : ﴿ دائرة(ن) ⋄ مستطيل(ع، ا) ﴾
مساحة ≡ λش. طابق ش : ﴿
    دائرة(ن) ⇒ 3 · ن · ن
    ⋄ مستطيل(ع، ا) ⇒ ع · ا
﴾
⎕ نص(مساحة(دائرة(5)))   # 75
\end{lstlisting}

\subsection{السمات}
\label{sec:design:traits}

\texttt{سمة قابلة\_للطباعة(ن) : ﴿ اطبع\_ذاتي : 0 ﴾} مع
\texttt{تطبيق قابلة\_للطباعة على نقطة : ﴿ اطبع\_ذاتي(ن) ≔ "نقطة" ﴾}. يُولَّد لكل
تطبيق closure داخلي باسم \texttt{\_\_trait\_\_{سمة}\_\_{نوع}\_\_{دالة}}،
ويُوجَّه الاستدعاء إليه بعد استنتاج نوع الوسيط من \texttt{type\_env}، مع
احتياطي (fallback) يحل الدالة في الدوال العامة \texttt{λس. احصل(س) · 2}.

\subsection{Result/Option ومعامل ؟}
\label{sec:design:result}

أربعة بناة مدمجة \texttt{نجاح(ق) | فشل(س) | بعض(ق) | لاشيء} بوسم موحد
(0/1/2/3). معامل \texttt{؟} ينفذ استخلاصاً مشروطاً (extract-or-propagate):

\begin{lstlisting}[style=arlang]
# Assembly المولَّد لمعامل ؟
mov rbx, [rax + 8]        # قراءة الوسم
test rbx, rbx
jnz .rq_fail_k            # إن فشل: اترك النتيجة الفاشلة
mov rax, [rax + 16]       # وإلا: استخرج القيمة
\end{lstlisting}

وهذا يُمكّن السلاسل الآمنة: \texttt{س١ ≔ اقسم(أ، ب)؟} تمرّر الفشل فوراً دون
تحقق صريح في كل نقطة.

\section{التقييم التجريبي}
\label{sec:evaluation}

\subsection{مجموعات الاختبار}
\label{sec:evaluation:suites}

\begin{table}[h]
\centering
\caption{نتائج مجموعات الاختبار}
\label{tab:suites}
\begin{tabular}{lcr}
\toprule
\textbf{السلسلة} & \textbf{عدد الاختبارات} & \textbf{النتيجة} \\
\midrule
test\_all\_phases.py & 24 & 24/24 \\
test\_phases.py (40–42) & 12 & 12/12 \\
test\_engine\_uas.py & 9 & 9/9 \\
phase43\_tests.py & 10 & 10/10 \\
test\_compound\_indexing.py & 9 & 9/9 \\
phase44\_adt\_tests.py & 9 & 9/9 \\
phase45\_traits\_tests.py & 9 & 9/9 \\
phase46\_result\_option\_tests.py & 9 & 9/9 \\
\midrule
\textbf{المجموع} & \textbf{91} & \textbf{91/91} \\
\bottomrule
\end{tabular}
\end{table}

\subsection{أمثلة مختارة}
\label{sec:evaluation:examples}

نمط عودي كامل:

\begin{lstlisting}[style=arlang]
مضروب ≡ λن. طابق ن : ﴿
    0 ⇒ 1
    ⋄ س ⇒ س · مضروب(س - 1)
﴾
⎕ نص(مضروب(5))   # 120
\end{lstlisting}

Result/Option متسلسل:

\begin{lstlisting}[style=arlang]
حساب ≡ λ(أ، ب، ج). طابق اقسم(أ، ب) : ﴿
    نجاح(ق١) ⇒ طابق اقسم(ق١، ج) : ﴿
        نجاح(ق٢) ⇒ نجاح(ق٢)
        ⋄ فشل(س) ⇒ فشل(س)
    ﴾
    ⋄ فشل(س) ⇒ فشل(س)
﴾
⎕ طابق حساب(100، 2، 5) : ﴿ نجاح(ق) ⇒ نص(ق) ⋄ فشل(س) ⇒ س ﴾   # 10
\end{lstlisting}

\subsection{الأداء}
\label{sec:evaluation:performance}

حجم الثنائي لبرنامج \emph{Hello World} حوالي 11KB، وبرنامج ADTs+Traits حوالي
12KB — لا يوجد runtime ولا startup overhead. وقت التجميع لبرنامج بسيط ~0.1s
ولبرنامج 1000 سطر ~2s (محدود بسرعة \texttt{nasm}، لا بالمولد). وقت التنفيذ
مقارب لـ C لأن الكود المولّد syscalls خام مباشرة.

\subsection{المقارنة}
\label{sec:evaluation:comparison}

\begin{table}[h]
\centering
\caption{المقارنة مع اللغات العالمية}
\label{tab:comparison}
\begin{tabular}{lccccc}
\toprule
\textbf{الميزة} & \textbf{لغتنا} & \textbf{Rust} & \textbf{Haskell} & \textbf{Go} \\
\midrule
AOT Compilation & ✔ & ✔ & ✔ & ✔ \\
No Runtime & ✔ & ✔ & ✘ (GC) & ✘ (GC) \\
Linear Ownership & ✔ & ✔ & ✘ & ✘ \\
Pattern Matching & ✔ & ✔ & ✔ & ✘ \\
Algebraic Data Types & ✔ & ✔ & ✔ & ✘ \\
Traits/Typeclasses & ✔ & ✔ & ✔ & ✘ \\
Arabic Syntax & ✔ & ✘ & ✘ & ✘ \\
No libc & ✔ & ✘ & ✘ & ✘ \\
\bottomrule
\end{tabular}
\end{table}

\section{العمل المستقبلي}
\label{sec:future}

الخارطة المتبقية: المرحلة 47 تضيف \texttt{بانتظار} (Async/Await) على حلقة
event loop مبنية على \texttt{epoll}، والمرحلة 48 تضيف Macros على نمط Rust
macro\_rules، والمرحلة 49 تستكشف Dependent Types (\texttt{نوع مصفوفة(ن، طول)}
حيث يُقيد الطول بالنوع)، والمرحلة 50 تستكمل self-hosting بحيث يجمّع المُجمّع
نفسه بدون Python نهائياً.

\section{الخلاصة}
\label{sec:conclusion}

قدمنا اللغة الرياضية العربية، أول لغة برمجة عربية كاملة تُجمَّع AOT إلى
ثنائيات ELF مستقلة. أثبتنا صحة التصميم بـ 91 اختباراً ناجحاً، وأعدنا إنتاج
أحدث تقنيات لغات البرمجة العالمية (الملكية الخطية، الأنواع الجبرية، السمات،
Result/Option) بصيغة عربية رياضية محكومة بدستور من ثمانية مبادئ. اللغة جاهزة
للاستخدام العملي وتشكل أساساً لمجتمع عربي للبرمجة.

\section*{الإهداء}

\begin{center}
\begin{minipage}{0.8\textwidth}
\centering
{\LARGE\bfseries بسم الله الرحمن الرحيم}

\medskip
هذا العمل إهداءٌ إلى:

\textbf{أمي ريحانة} — قلبي وروحه.

\medskip
حمداني سيدي محمد

HAMDANI SIDI MOHAMED

﴿وقل رب زدني علماً﴾
\end{minipage}
\end{center}

\bibliographystyle{ACM-Reference-Format}
\bibliography{references}

\end{document}

